% Fragmento integrado desde main.tex con % Fragmento integrado desde main.tex con % Fragmento integrado desde main.tex con % Fragmento integrado desde main.tex con \input{citas_oficiales.tex}.
% Para BibTeX clásico, agregar al final del informe:
% \bibliographystyle{IEEEtran}
% \bibliography{referencias}
%
% Con biblatex, usar en el preámbulo:
% \addbibresource{referencias.bib}
% y \printbibliography en lugar de los dos comandos anteriores.

\subsection{Flujo RTL a GDSII con LibreLane}
La implementación ASIC se ejecutó con LibreLane 3.0.16 en Docker/WSL. LibreLane
acepta configuraciones JSON validadas y organiza un flujo reproducible de
implementación digital desde RTL hasta las vistas físicas. La configuración
reglamentaria usa \texttt{DESIGN\_NAME=asic\_top}, \texttt{CLOCK\_PORT=clk\_i},
\texttt{CLOCK\_PERIOD=100 ns} y \texttt{FP\_CORE\_UTIL=40}.
La documentación oficial se usa para describir el flujo y sus parámetros;
las métricas de este proyecto provienen exclusivamente de sus logs y reportes
originales~\cite{librelane_docs}.

\subsection{Tecnología GF180MCU y biblioteca estándar}
El flujo utiliza la variante instalada \texttt{gf180mcuD}, correspondiente a
la familia de tecnología GF180MCU, con la revisión de \texttt{open\_pdks}
\texttt{8afc8346a57fe1ab7934ba5a6056ea8b43078e71}. La biblioteca seleccionada
es \texttt{gf180mcu\_fd\_sc\_mcu7t5v0}, una biblioteca estándar de siete
tracks documentada por el PDK. Esta referencia respalda la procedencia de las
vistas tecnológicas y de celdas, no reemplaza la revisión exacta registrada en
las corridas~\cite{gf180mcu_pdk_docs}.

\subsection{Implementación física y análisis de timing con OpenROAD}
LibreLane invoca OpenROAD para inicialización del floorplan, placement, CTS,
routing, extracción RC y STA post-route. El uso de utilización para dimensionar
el core, la inspección de congestión y el análisis después de extracción son
coherentes con las capacidades documentadas de OpenROAD~\cite{openroad_docs}.
Los valores de WNS, TNS, DRC, LVS, antena y violaciones eléctricas presentados
en este informe se toman de los reportes originales de LibreLane/OpenROAD; no
se deducen de esta documentación.
.
% Para BibTeX clásico, agregar al final del informe:
% \bibliographystyle{IEEEtran}
% \bibliography{referencias}
%
% Con biblatex, usar en el preámbulo:
% \addbibresource{referencias.bib}
% y \printbibliography en lugar de los dos comandos anteriores.

\subsection{Flujo RTL a GDSII con LibreLane}
La implementación ASIC se ejecutó con LibreLane 3.0.16 en Docker/WSL. LibreLane
acepta configuraciones JSON validadas y organiza un flujo reproducible de
implementación digital desde RTL hasta las vistas físicas. La configuración
reglamentaria usa \texttt{DESIGN\_NAME=asic\_top}, \texttt{CLOCK\_PORT=clk\_i},
\texttt{CLOCK\_PERIOD=100 ns} y \texttt{FP\_CORE\_UTIL=40}.
La documentación oficial se usa para describir el flujo y sus parámetros;
las métricas de este proyecto provienen exclusivamente de sus logs y reportes
originales~\cite{librelane_docs}.

\subsection{Tecnología GF180MCU y biblioteca estándar}
El flujo utiliza la variante instalada \texttt{gf180mcuD}, correspondiente a
la familia de tecnología GF180MCU, con la revisión de \texttt{open\_pdks}
\texttt{8afc8346a57fe1ab7934ba5a6056ea8b43078e71}. La biblioteca seleccionada
es \texttt{gf180mcu\_fd\_sc\_mcu7t5v0}, una biblioteca estándar de siete
tracks documentada por el PDK. Esta referencia respalda la procedencia de las
vistas tecnológicas y de celdas, no reemplaza la revisión exacta registrada en
las corridas~\cite{gf180mcu_pdk_docs}.

\subsection{Implementación física y análisis de timing con OpenROAD}
LibreLane invoca OpenROAD para inicialización del floorplan, placement, CTS,
routing, extracción RC y STA post-route. El uso de utilización para dimensionar
el core, la inspección de congestión y el análisis después de extracción son
coherentes con las capacidades documentadas de OpenROAD~\cite{openroad_docs}.
Los valores de WNS, TNS, DRC, LVS, antena y violaciones eléctricas presentados
en este informe se toman de los reportes originales de LibreLane/OpenROAD; no
se deducen de esta documentación.
.
% Para BibTeX clásico, agregar al final del informe:
% \bibliographystyle{IEEEtran}
% \bibliography{referencias}
%
% Con biblatex, usar en el preámbulo:
% \addbibresource{referencias.bib}
% y \printbibliography en lugar de los dos comandos anteriores.

\subsection{Flujo RTL a GDSII con LibreLane}
La implementación ASIC se ejecutó con LibreLane 3.0.16 en Docker/WSL. LibreLane
acepta configuraciones JSON validadas y organiza un flujo reproducible de
implementación digital desde RTL hasta las vistas físicas. La configuración
reglamentaria usa \texttt{DESIGN\_NAME=asic\_top}, \texttt{CLOCK\_PORT=clk\_i},
\texttt{CLOCK\_PERIOD=100 ns} y \texttt{FP\_CORE\_UTIL=40}.
La documentación oficial se usa para describir el flujo y sus parámetros;
las métricas de este proyecto provienen exclusivamente de sus logs y reportes
originales~\cite{librelane_docs}.

\subsection{Tecnología GF180MCU y biblioteca estándar}
El flujo utiliza la variante instalada \texttt{gf180mcuD}, correspondiente a
la familia de tecnología GF180MCU, con la revisión de \texttt{open\_pdks}
\texttt{8afc8346a57fe1ab7934ba5a6056ea8b43078e71}. La biblioteca seleccionada
es \texttt{gf180mcu\_fd\_sc\_mcu7t5v0}, una biblioteca estándar de siete
tracks documentada por el PDK. Esta referencia respalda la procedencia de las
vistas tecnológicas y de celdas, no reemplaza la revisión exacta registrada en
las corridas~\cite{gf180mcu_pdk_docs}.

\subsection{Implementación física y análisis de timing con OpenROAD}
LibreLane invoca OpenROAD para inicialización del floorplan, placement, CTS,
routing, extracción RC y STA post-route. El uso de utilización para dimensionar
el core, la inspección de congestión y el análisis después de extracción son
coherentes con las capacidades documentadas de OpenROAD~\cite{openroad_docs}.
Los valores de WNS, TNS, DRC, LVS, antena y violaciones eléctricas presentados
en este informe se toman de los reportes originales de LibreLane/OpenROAD; no
se deducen de esta documentación.
.
% Para BibTeX clásico, agregar al final del informe:
% \bibliographystyle{IEEEtran}
% \bibliography{referencias}
%
% Con biblatex, usar en el preámbulo:
% \addbibresource{referencias.bib}
% y \printbibliography en lugar de los dos comandos anteriores.

\subsection{Flujo RTL a GDSII con LibreLane}
La implementación ASIC se ejecutó con LibreLane 3.0.16 en Docker/WSL. LibreLane
acepta configuraciones JSON validadas y organiza un flujo reproducible de
implementación digital desde RTL hasta las vistas físicas. La configuración
reglamentaria usa \texttt{DESIGN\_NAME=asic\_top}, \texttt{CLOCK\_PORT=clk\_i},
\texttt{CLOCK\_PERIOD=100 ns} y \texttt{FP\_CORE\_UTIL=40}.
La documentación oficial se usa para describir el flujo y sus parámetros;
las métricas de este proyecto provienen exclusivamente de sus logs y reportes
originales~\cite{librelane_docs}.

\subsection{Tecnología GF180MCU y biblioteca estándar}
El flujo utiliza la variante instalada \texttt{gf180mcuD}, correspondiente a
la familia de tecnología GF180MCU, con la revisión de \texttt{open\_pdks}
\texttt{8afc8346a57fe1ab7934ba5a6056ea8b43078e71}. La biblioteca seleccionada
es \texttt{gf180mcu\_fd\_sc\_mcu7t5v0}, una biblioteca estándar de siete
tracks documentada por el PDK. Esta referencia respalda la procedencia de las
vistas tecnológicas y de celdas, no reemplaza la revisión exacta registrada en
las corridas~\cite{gf180mcu_pdk_docs}.

\subsection{Implementación física y análisis de timing con OpenROAD}
LibreLane invoca OpenROAD para inicialización del floorplan, placement, CTS,
routing, extracción RC y STA post-route. El uso de utilización para dimensionar
el core, la inspección de congestión y el análisis después de extracción son
coherentes con las capacidades documentadas de OpenROAD~\cite{openroad_docs}.
Los valores de WNS, TNS, DRC, LVS, antena y violaciones eléctricas presentados
en este informe se toman de los reportes originales de LibreLane/OpenROAD; no
se deducen de esta documentación.
